\section{Future Scope}
\label{sec:future}

\subsection{Status of the items carried over from Semester VI}

Table~\ref{tab:carryover} revisits every direction proposed in the previous report and records what
the present phase did with it.

\begin{table}[H]
\centering
\caption{Semester VI proposals and their current status.}
\label{tab:carryover}
\footnotesize
\begin{tabular}{@{}L{4.4cm}L{2.2cm}L{8.3cm}@{}}
\toprule
\textbf{Proposal} & \textbf{Status} & \textbf{Position at the end of Semester VII} \\
\midrule
Examination scheduling & still open &
  Examination periods are declared in \texttt{1sessionSetup.xml}, but no per-course examination
  elements are emitted and the examination module has not been exercised. Taasika does not model
  examinations, so the data would have to come from elsewhere. \\
Calendar (\texttt{.ics}) integration & still open &
  Not started. The exported CSV contains everything a calendar export needs, so this is now a
  self-contained task rather than one blocked on the timetable. \\
University-wide deployment & still open &
  The pipeline is parameterised by snapshot and term and is not specific to one department, but
  only Computer Engineering data has been loaded and solved. \\
Automated change notification & still open & Not started. \\
WhatsApp group automation & still open & Not started. \\
Student data pipeline from the MIS portal & partly done &
  The transformation and import path is built, validated and exercised end to end, producing
  1{,}130 student records with 9{,}148 enrolments. Only the data source is missing: the generator
  simulates students because MIS extracts were not available. Substituting the real source is a
  change to one generator. \\
Course definition interface for instructors & partly done &
  Course, section and preference generation is fully automated from the Taasika snapshot, so no
  manual data entry is required. The instructor-facing interface for submitting availability has
  not been built. \\
Constraint mapping onto UniTime & done &
  Established in Semester VI and now implemented as generated data rather than as a mapping on
  paper (Section~\ref{sec:constraints}). \\
Handling of shared multi-instructor laboratories & partly done &
  Multiple lead instructors on one class is supported by UniTime and used, but the stronger
  requirement---two instructors blocked simultaneously for the same student set in the same
  room---has no distribution constraints generated for it. \\
\bottomrule
\end{tabular}
\end{table}

\subsection{New directions arising from this phase}

\begin{enumerate}
  \item \textbf{Resolve the residual first-year conflicts.} The 40 remaining conflicts are
    concentrated in four laboratory placements against the reserved first-year lecture block.
    Adding generated different-time distribution constraints between the programming laboratory
    subpart and the reserved block would make the requirement explicit to the solver rather than
    leaving it to the student-conflict weight. This is the smallest remaining step with a direct
    effect on the headline metric.

  \item \textbf{Student sectioning with pre-defined batches.} The institution assigns batches
    administratively, which is why student switching had to be disabled. The principled alternative
    is to express the batches as group reservations, so that UniTime may section students but only
    within the batch they belong to. This would recover the solver's freedom to resolve conflicts
    by sectioning without discarding the cohort structure.

  \item \textbf{Soft-constraint optimisation.} The room preference hit rate of 12.59\,\% is the
    weakest measured quantity. Work is needed on both sides: tightening the generated preferences,
    which are currently permissive, and tuning the room-preference weight against the
    student-conflict weight in the manner of the sweep in Section~\ref{sec:experiments}.

  \item \textbf{Fix and extend the evaluator.} Normalising instructor names would repair the
    identity check and make the preference score meaningful. Replacing the exact-slot matching in
    \texttt{sectioning\_slot\_conflicts.py} with the interval test used elsewhere would remove a
    tool that currently reports false negatives.

  \item \textbf{One-command snapshot to solution.} Generating, importing, solving and evaluating
    currently requires three scripts and several manual steps in the UniTime interface. Since the
    headless solver can already load a solver XML and solve it outside the server, an end-to-end
    command from a Taasika snapshot to a scored solution is achievable and would make the
    refinement loop of Section~\ref{sec:evaluation} substantially faster.

  \item \textbf{Write the solution back to Taasika.} The timetable is currently a dead end: it is
    exported as CSV and XML but never returned to the system the institution actually uses. A
    converter from a solved timetable into the Taasika snapshot format would close the loop and is
    a precondition for the system being adopted operationally.

  \item \textbf{Complete the learning-layer evaluation, or retire it.} Arms L and R should be run
    against arm B1 on a held-out term, at equal wall-clock budget, on an instance that is not
    already conflict-free. If a harder instance---a university-wide one, for example---proves
    resistant to weight tuning, the layer is ready to be evaluated on it. If no such instance
    appears, the honest conclusion of Section~\ref{sec:rl} stands and the layer should be kept as
    instrumentation rather than as a solver component.

  \item \textbf{Correct the dataset defects.} The unplaced first-year cohort, the over-subscribed
    M.Tech laboratory, the cross-division batches, the empty groups and the obsolete laboratory
    sections 16--41 are all input-data problems with known fixes, and each removes a source of
    spurious conflicts.
\end{enumerate}
